\section{The minimizing branches and the predictor derivative}
\label{sec:predictor-branch}
We establish the branch properties used in the local comparison.
The cutoffs for uniqueness can be chosen using only the group, before
choosing a patch. Analyticity near the identity is then a local
finite-dimensional statement, and its derivative includes the root
normalization explicitly.

\subsection{The nearest-point projection onto the group}
The gauge group $G$ is a closed subgroup of $U(N)$. We regard
$\mathcal M=M_N(\mathbb C)$ as a real vector space with the Frobenius
inner product $\langle A,B\rangle=\operatorname{Re}\operatorname{tr}(A^*B)$
and norm $\|\cdot\|$. Left and right multiplication by a unitary
matrix are isometries of $\mathcal M$. Let
$\mathfrak g=\{X\in\mathcal M:\exp(sX)\in G\text{ for all }s\in\mathbb R\}$
be the Lie algebra of $G$, a real subspace of the skew-Hermitian
matrices, let $\mathfrak g^\perp$ be its orthogonal complement in
$\mathcal M$, and let $\Pi$ be the orthogonal projection of $\mathcal M$
onto $\mathfrak g$. A point $u\in G$ is \emph{nearest} to
$Z\in\mathcal M$ if $\|Z-u\|\le\|Z-v\|$ for all $v\in G$. Since $G$ is
compact, every $Z$ has at least one nearest point. The group may be
disconnected, and it may be finite, in which case $\mathfrak g=0$.

\begin{lemma}\label{lem:branch-nearest-projection}
There are $\rho>0$ and a real-analytic map
$\pi:B(1,2\rho)\to G$, defined on the open Frobenius ball of radius
$2\rho$ about $1$ in $\mathcal M$, with the following properties.
\begin{enumerate}
\item[(a)] $\pi(u)=u$ for $u\in G\cap B(1,2\rho)$, and $D\pi(1)=\Pi$.
\item[(b)] For every $Z\in B(1,\rho)$, $\pi(Z)$ is the only point of $G$
  nearest to $Z$.
\item[(c)] If $Z\in\mathcal M$ and $\|Z-v\|<\rho$ for some $v\in G$, then
  $Z$ has exactly one nearest point in $G$, namely $v\pi(v^{-1}Z)$.
\item[(d)] The first and second derivatives of $\pi$ are bounded on the
  closed ball $\overline B(1,\rho)$.
\end{enumerate}
\end{lemma}

\begin{proof}
We use one fact from the theory of matrix Lie groups. Since $U(N)$ is
closed in $GL(N,\mathbb C)$, $G$ is a closed subgroup of
$GL(N,\mathbb C)$, that is a matrix Lie group in the sense of B.~C. Hall,
\emph{Lie Groups, Lie Algebras, and Representations}, 2nd ed.,
Graduate Texts in Mathematics 222, Springer, 2015, Definition~1.4,
and $\mathfrak g$ is its Lie algebra in the sense of Definition~3.18
there. By Theorem~3.42 of the same book there is $\varepsilon_0\in(0,\log2)$
such that, for every matrix $A=\exp X$ with $\|X\|<\varepsilon_0$,
one has $A\in G$ if and only if $X\in\mathfrak g$. The matrix
exponential is a real-analytic diffeomorphism from a neighbourhood
of $0$ in $\mathcal M$ onto a neighbourhood of $1$, with derivative
the identity at $0$, so there is $\theta>0$ such that every
$u\in G$ with $\|u-1\|<\theta$ equals $\exp X$ for one
$X\in\mathfrak g$ with $\|X\|<\varepsilon_0$, and $X\to0$ as
$u\to1$. (Equivalently: near $1$, $G$ is the embedded analytic
submanifold $\exp(\mathfrak g)$; this is Cartan's closed-subgroup
theorem.)

\emph{The normal chart.} Define
$\Phi:\mathfrak g\times\mathfrak g^\perp\to\mathcal M$ by
$\Phi(X,N)=\exp(X)(1+N)$. It is real-analytic, $\Phi(0,0)=1$, and
$D\Phi(0,0)(X,N)=X+N$. Since $\mathcal M=\mathfrak g\oplus\mathfrak g^\perp$,
this derivative is a linear isomorphism. By the real-analytic inverse
function theorem (S.~G. Krantz and H.~R. Parks, \emph{A Primer of Real
Analytic Functions}, 2nd ed., Birkh\"auser, 2002, Section~2.5) there is
$r\in(0,\varepsilon_0)$ such that, with
$O_r=\{(X,N):\|X\|<r,\ \|N\|<r\}$, the set $\Phi(O_r)$ is an open
neighbourhood of $1$ and $\Phi:O_r\to\Phi(O_r)$ is a real-analytic
diffeomorphism. Choose $\rho>0$ so small that
$\overline B(1,3\rho)\subset\Phi(O_r)$, $2\rho<\theta$, $\rho\le r$, and
every $u\in G$ with $\|u-1\|<2\rho$ is $\exp X$ with
$X\in\mathfrak g$, $\|X\|<r$. For $Z\in B(1,2\rho)$ write
$\Phi^{-1}(Z)=(X(Z),N(Z))$ and put $\pi(Z)=\exp X(Z)$. Then $\pi$ is
real-analytic and takes values in $G$, because $X(Z)\in\mathfrak g$.

\emph{Property (a).} If $u\in G\cap B(1,2\rho)$, then $u=\exp X$ with
$X\in\mathfrak g$, $\|X\|<r$, so $u=\Phi(X,0)$ with $(X,0)\in O_r$.
Injectivity of $\Phi$ on $O_r$ gives $X(u)=X$ and $\pi(u)=u$.
Differentiating $\pi(\Phi(X,N))=\exp X$ at $(0,0)$ gives
$D\pi(1)(X+N)=X$, that is $D\pi(1)=\Pi$.

\emph{Property (b).} Let $Z\in B(1,\rho)$ and let $u$ be a nearest
point. Since $1\in G$, $\|Z-u\|\le\|Z-1\|<\rho$ and therefore
$\|u-1\|<2\rho$, so $u=\exp X$ with $X\in\mathfrak g$, $\|X\|<r$.
For $Y\in\mathfrak g$, $s\mapsto\|Z-u\exp(sY)\|^2$ is differentiable
and minimal at $s=0$, because $u\exp(sY)\in G$. Its derivative there
is $-2\langle Z-u,uY\rangle=-2\langle u^*Z-1,Y\rangle$, using that $u$
is unitary. Hence $N:=u^*Z-1\in\mathfrak g^\perp$, and
$\|N\|=\|Z-u\|<\rho\le r$. Thus $Z=\Phi(X,N)$ with $(X,N)\in O_r$, so
$(X,N)=\Phi^{-1}(Z)$ and $u=\pi(Z)$. As a nearest point exists,
$\pi(Z)$ is the unique one.

\emph{Property (c).} Left multiplication by $v^{-1}\in G$ is an
isometry of $\mathcal M$ mapping $G$ onto $G$. Hence $u$ is nearest to
$Z$ if and only if $v^{-1}u$ is nearest to $v^{-1}Z$, and
$\|v^{-1}Z-1\|=\|Z-v\|<\rho$. By (b), $v^{-1}u=\pi(v^{-1}Z)$.

\emph{Property (d).} The derivatives of $\pi$ are continuous on the
compact set $\overline B(1,\rho)\subset B(1,2\rho)$.
\end{proof}

In the rest of this chapter and in Sections~\ref{sec:seed-geometry}
and~\ref{sec:seed-energy}, $\pi$ denotes this map. Its construction
uses only $G$; no other part of the data enters.

\subsection{Uniform uniqueness of minimizing frames}
Use the finite path family $W_{ij\alpha}$ and minimizing energy from
Section~\ref{sec:seed-geometry}. Let $\pi$ and $\rho$ be as in
Lemma~\ref{lem:branch-nearest-projection}. By continuity of
$D\pi$ at $1$ and $D\pi(1)=\Pi$, there is a convex ball about $1$,
contained in $B(1,\rho)$, on which $\|D\pi-\Pi\|\le\tfrac14$; for all
$A,B$ in that ball,
\begin{equation}\label{eq:branch-projection}
 \|\pi(A)-\pi(B)-\Pi(A-B)\|\le\tfrac14\|A-B\|.
\end{equation}
This follows by integrating $D\pi-\Pi$ along the segment. Choose
$\delta>0$ for such a ball, and choose
$0<\varepsilon\le\min\{\delta/4,1/4\}$. Suppose
\[
 \|W_{ij\alpha}-1\|\le\varepsilon
 \quad\text{for every }i,j,\alpha.
\]
We claim that any two minimizing frames differ by a common left
multiplication, independently of the number of points or paths.

Comparison with the constant frame and the squared triangle
inequality give
\[
 \mathop{\rm av}_{i,j}\|h_i-h_j\|^2\le4\varepsilon^2.
\]
For clarity, the mean squared transported discrepancy
$\|h_iW_{ij\alpha}-h_j\|^2$ is at most $\varepsilon^2$ by
minimality. Each untransported discrepancy squared is at most twice
this term plus $2\varepsilon^2$; averaging gives the claim.
Rotate each minimizing frame separately so that $1$ is nearest to
its centroid $m$. Choosing a point with mean squared distance no
larger than the double average and applying Jensen gives
$\|m-1\|\le2\varepsilon$. Its feedback matrix
$m_i=\mathop{\rm av}_{j,\alpha}h_jW_{ij\alpha}^{-1}$ satisfies
$\|m_i-m\|\le\varepsilon$, so $\|m_i-1\|\le3\varepsilon<\delta$.
By \eqref{eq:geometry-feedback}, $h_i$ is a nearest group element to
$m_i$, and $m_i$ lies in the $\delta$-ball, which is contained in
$B(1,\rho)$. Hence
Lemma~\ref{lem:branch-nearest-projection}(b) identifies $h_i=\pi(m_i)$.

Let $h,k$ be the two rotated minimizers and set
$v=\max_i\|h_i-k_i\|$. Their feedback differences $D_i$ satisfy
\[
 \|D_i\|\le v,\qquad
 \|D_i-(m_h-m_k)\|\le v\varepsilon.
\]
The first follows by averaging and unitary invariance. For the second,
subtract the centroid difference to obtain the average of
$(h_j-k_j)(W_{ij\alpha}^{-1}-1)$ and use the matrix product bound.
Since $1$ is nearest to each centroid,
$\Pi(m_h-1)=\Pi(m_k-1)=0$: differentiate squared distance along
$\exp(tX)$ for each tangent vector $X$. Thus
$\|\Pi D_i\|\le v\varepsilon$. At an index attaining the maximum,
\eqref{eq:branch-projection} yields
\[
 v\le v\varepsilon+\tfrac14v\le\tfrac12v,
\]
so $v=0$. Undoing the rotations proves uniqueness up to a common
left multiplier. If both frames are normalized at their centre,
that multiplier is $1$.

The actual frame branch tests that, for every rooted axial gauge
centred at the patch centre (every coordinate order of the axial
tree), every coordinate path transport has chord less than a fixed
cutoff. Take this cutoff at most $\varepsilon/D$. Choose one of
those axial gauges; the resulting $W$ satisfies the preceding matrix
bound. Changing variables by that gauge in the minimizing energy
proves uniqueness of the normalized frame on the entire tested
branch. These cutoffs do not depend on the patch size or smoothing
order. At the identity all tested chords vanish, so this branch is
an open neighbourhood of the identity field.

For the predictor branch, let $w_i\ge0$, $\sum_iw_i=1$, and let
$m=\sum_iw_iY_i$. Expanding the squares shows that minimizing
$\sum_iw_i r(Y_iu^{-1})^2/2$ is equivalent to choosing a nearest
group element to $m$. If the weighted cost at some $u$ is less than
$\zeta$, Jensen gives
\[
 \|m-u\|^2\le\sum_iw_i\|Y_i-u\|^2<2D^2\zeta.
\]
Choose $\zeta$ so that $\sqrt{2D^2\zeta}<\rho$. Then $\|m-u\|<\rho$
with $u\in G$, and Lemma~\ref{lem:branch-nearest-projection}(c),
the uniform nearest-point tube obtained by left translation, shows
that $m$ has exactly one nearest group element. Then the predictor
minimum is unique on its tested branch, uniformly over all finite
weight distributions. Finally the fallback frame rule has an
eligible choice by \eqref{eq:geometry-axial} with axial factor $16$.
Together these facts construct the required branch constants for any
fixed positive predictor penalty.

\subsection{Analyticity of the normalized frame}
Fix a patch and a finite torus for the moment. Write $a$ for the
fine logarithmic bond field and $h_i(a)$ for the actual frame
matrices. At $a=0$ they are all $1$. The chosen minimizer is
continuous on the uniqueness branch. To verify this without assuming
regularity of a choice function, consider a convergent sequence of
input fields. Compactness of the finite product of groups gives a
convergent subsequence of any sequence of minimizers. Continuity of
the energy and the centre constraint makes every limit a minimizer
at the limiting field. Uniqueness makes that limit the selected
one. A sequence staying outside a fixed neighbourhood of that value
would contradict this argument, which proves continuity.

Let $o$ be the patch centre and let $W_{ij\alpha}(a)$ be its actual
coordinate path transports. Define an ambient matrix-valued system
\[
 \mathcal F_i(a,H)=
 \begin{cases}
 H_o-1,&i=o,\\
 H_i-\pi\bigl(\mathop{\rm av}_{j,\alpha}
                     H_jW_{ij\alpha}(a)^{-1}\bigr),&i\ne o.
 \end{cases}
\]
It is analytic near $(0,\mathbf1)$. Continuity of the actual frame
and its feedback, together with self-consistency, shows that
$\mathcal F(a,h(a))=0$ near zero. The derivative in the ambient
frame variable is
\[
 (\mathcal D v)_o=v_o,\qquad
 (\mathcal D v)_i=v_i-\Pi\bar v\quad(i\ne o),
 \qquad \bar v=\mathop{\rm av}_jv_j.
\]
It is invertible. Indeed if $\mathcal Dv=0$, put $A=\Pi\bar v$.
Then $v_o=0$ and $v_i=A$ for $i\ne o$. If the patch has $n$ points,
$\bar v=(n-1)A/n$, so applying $\Pi$ gives $A=(n-1)A/n$ and
$A=0$. The derivative is an injective endomorphism of a finite
-dimensional space and is therefore bijective. The analytic implicit
function theorem gives a unique analytic local solution. Continuity
places the actual frame in that local solution neighbourhood, so
$h(a)$ is analytic. This argument works in ambient matrices; it does
not require choosing coordinates for a complement at every frame.

\subsection{The root correction in the derivative}
For a Lie-algebra-valued bond perturbation $a$, write
$\int_\gamma a$ for its signed sum along a path. Define
\[
 \tau_y(i;a)=-\mathop{\rm av}_{j,\alpha}
                 \int_{i\to j,\alpha}a,
\]
where the paths are in the patch centred at the coarse site $y$.
Throughout, a coarse site $y$ is embedded in the fine lattice as $Ly$,
and the patch at $y$ consists of the fine points $Ly+i$, with $i$
ranging over the patch offsets in fine units and periodic reduction
understood; thus $\int_{i\to j,\alpha}a$ is the sum of $a$ along the
path from $Ly+i$ to $Ly+j$ of order $\alpha$.
Differentiate the feedback equation at zero. If
$\dot h_i$ is the matrix derivative and $\overline{\dot h}$ its
uniform average, differentiation of the inverse transport gives
\[
 \dot h_i=\Pi\overline{\dot h}+\operatorname{col}(\tau_y(i;a)).
\]
The feedback identity holds also at the centre: the cost is
invariant under a common left multiplication, so the
centre-normalized minimizer is an unconstrained minimizer and
\eqref{eq:geometry-feedback} applies at $o$ as well (see
Section~\ref{sec:seed-geometry}). Since the normalized
frame has $h_o(a)=1$, its centre derivative is zero. Subtraction
therefore yields
\begin{equation}\label{eq:branch-frame-derivative}
 \dot h_i=\operatorname{col}\{\tau_y(i;a)-\tau_y(o;a)\}.
\end{equation}
In particular the frame derivative uses uniform averages over patch
points and path orders, before the separate connector weights enter.

For a coarse bond $e=(y,b)$, the connector at patch point $i$ consists
of its frame at $y$, the length-$L$ fine path, and the inverse frame
at $y+e_b$. Its colour derivative is consequently
\[
 \tau_y(i;a)-\tau_y(o;a)
 +\int_{Ly+i\to L(y+e_b)+i}a
 -\tau_{y+e_b}(i;a)+\tau_{y+e_b}(o;a).
\]
At zero all connectors equal $1$. Nearby the predictor is on its
minimizing branch and equals $\pi$ of their weighted matrix centroid.
This follows from the weighted nearest-point identity above; the
branch remains open since its minimum cost is at most the continuous
cost of comparison with $1$, which vanishes at zero. Thus the
predictor is analytic near zero. The derivative of $\pi$ is $\Pi$,
which fixes the above colour derivatives, so the predictor derivative
is their weighted sum.

Define the bond averaging map and the site correction by
\[
 \begin{split}
 (\mathcal Aa)_{y,b}
   &=\sum_iw_i\int_{Ly+i\to L(y+e_b)+i}a,\\
 (\mathcal Ra)_y&=\tau_y(o;a)-\sum_iw_i\tau_y(i;a).
 \end{split}
\]
The weighted connector derivative is exactly
\begin{equation}\label{eq:branch-probe}
 \mathcal Qa=\mathcal Aa+d_c^{(0)}\mathcal Ra.
\end{equation}
The sign is determined by the subtraction of the frame at the final
endpoint: the remaining root terms are
$(\mathcal Ra)_{y+e_b}-(\mathcal Ra)_y$.
For a fine gradient $a=d_f^{(0)}v$, path sums telescope, giving
\[
 (\mathcal R d_f^{(0)}v)_y
       =v(Ly)-\sum_iw_i v(Ly+i),\qquad
 \mathcal Qd_f^{(0)}v=d_c^{(0)}\{v(Ly)\}_y.
\]
This checks the root normalization directly. Equations
\eqref{eq:branch-frame-derivative} and \eqref{eq:branch-probe} prove
the analytic predictor law and derivative used in
Section~\ref{sec:seed-unhit-comparison}. Its transplantation argument
then supplies a uniform fixed-patch Taylor radius for all periods.

The covariance needed there follows from the same definitions.
A gauge change turns a frame into
$f'_i=g_o f_i g_i^{-1}$. On a minimizing branch it preserves the
energy and centre normalization, so uniqueness identifies this with
the selected frame; on a fallback branch the identity follows by
telescoping its rooted paths. Connector endpoint factors then reduce
to $g_{Ly}Y_i g_{L(y+e_b)}^{-1}$. The weighted cost is invariant under
this change of variables. Uniqueness identifies the minimizing
predictor, and direct substitution handles a selected connector on
the other branch. The original tags are retained in a gauge change.
These identities give the block-constant covariance used in the
unhit comparison.
